\section{Background --- draft v1}
\label{sec:background}

\subsection{Vocabulary}

Four terms carry precise meanings in this paper. A \textbf{diagnostic} is
a structured value describing one user-source mistake: kind,
severity, source span, message, root cause, explanation, ordered
fixes, optional invalid/valid examples, and a trace. A \textbf{trace} is
the sequence of compiler activities that led to a diagnostic's
creation; this paper distinguishes \emph{recorded} traces (snapshots of
the frames actually open at creation time) from \emph{curated} ones
(hand-written reference chains --- the fallback, \S6.3). \textbf{Provenance}
is the backward link from a compiled artifact to its origin: which
source line an IR instruction or assembly block derives from, and
which optimization decisions transformed it. \textbf{Explanation} is the
umbrella for all of the above --- preserved compiler evidence, not
machine-learning explainability and not post-hoc rationalization.

\subsection{The source language}

The vehicle is a C-like kernel language --- deliberately not "a C++
subset" in any meaningful sense, though its programs are valid C.
The complete grammar (verbatim from the parser's specification
comment):

\begin{verbatim}
program       -> function_decl*  EOF
function_decl -> type IDENT '(' param_list ')' block
param_list    -> (type IDENT (',' type IDENT)*)  |  (empty)
block         -> '{' statement* '}'
statement     -> var_decl | return_stmt
var_decl      -> type IDENT ('=' expression)? ';'
return_stmt   -> 'return' expression? ';'
expression    -> assignment
assignment    -> IDENTIFIER '=' assignment | equality
equality      -> comparison (('=='|'!=') comparison)*
comparison    -> addition   (('<'|'>') addition)*
addition      -> multiply   (('+'|'-') multiply)*
multiply      -> unary      (('*'|'/') unary)*
unary         -> primary
primary       -> INTEGER_LITERAL | IDENTIFIER | '(' expression ')'
type          -> 'int'
\end{verbatim}

Table T1 (feature inventory): \textbf{supported} --- \texttt{int} variables with
optional initializers; arithmetic and comparison operators with C
precedence; assignment as a right-associative expression
(\texttt{a = b = c}); multi-parameter functions; multi-error reporting with
panic-mode recovery. \textbf{Absent by design} --- control flow, function
calls, strings, arrays, pointers, unary operators, additional types.
Internally the type lattice contains \texttt{int}/\texttt{bool}/\texttt{unknown} (\texttt{bool}
arises only as a comparison result; \texttt{unknown} is the error-recovery
type that suppresses cascading type errors) and \texttt{void} (reserved).

The smallness is the method, not a concession: every architectural
mechanism must be visible whole, every diagnostic kind enumerable
(seventeen, of which five are explicitly reserved-unreachable), and
every claim checkable against a codebase a reviewer can actually
read (\textasciitilde{}5.9K lines, \S8).

\subsection{The pipeline}

Compilation proceeds through eight stages: lexing (single-pass,
maximal-munch), recursive-descent parsing to an AST whose nodes are
pure data traversed by four independent visitors, semantic analysis
(scope-stack symbol table; type resolution annotating the AST),
lowering to a flat three-address IR, three optimization passes
(constant folding, copy propagation, dead-code elimination) run
together to a fixed point, code generation to 32-bit x86 assembly
text, and an OS-facing driver that assembles and links a real
executable via the system toolchain. The construction follows
standard practice throughout \cite{aho2006compilers}; nothing in the
pipeline itself is claimed as a contribution. What \S6 adds is the
explanation obligation threaded through every one of these stages ---
and \S9 the bill for it.



